% Generated draft: requires mathematical and source review.
\documentclass[11pt]{article}
\usepackage[margin=25mm]{geometry}
\usepackage{fontspec}
\setmainfont{Latin Modern Roman}
\usepackage{amsmath,amssymb,mathtools,unicode-math}
\setmathfont{Latin Modern Math}
\usepackage{enumitem,needspace,xurl,hyperref}
\setlength{\parindent}{0pt}
\setlength{\parskip}{5pt}
\setlength{\emergencystretch}{3em}
\newcommand{\N}{\mathbb{N}}
\newcommand{\Z}{\mathbb{Z}}
\newcommand{\Q}{\mathbb{Q}}
\newcommand{\R}{\mathbb{R}}
\newcommand{\C}{\mathbb{C}}
\newcommand{\F}{\mathbb{F}}
\newcommand{\A}{\mathbb{A}}
\newcommand{\PP}{\mathbb{P}}
\newcommand{\E}{\mathbb{E}}
\newcommand{\eps}{\varepsilon}
\newcommand{\dd}{\,\mathrm{d}}
\begin{document}
\section*{UnsolvedMath ID: 11000022}
\textbf{Source problem code:} catalogue entry
\subsection{Definitions and mathematical statement}
Fix the genus \(g\), the source's surface convention, and a finite generating set \(S\) of \(\operatorname{Mod}_g\). The word problem takes a word \(w\in(S\cup S^{-1})^*\), of length \(n\), and decides whether its represented mapping class is the identity.

\emph{Question 3.1 (fast word problem).} Is there a subquadratic-time algorithm for this decision problem? With a computational model fixed, a weak precise interpretation asks for worst-case time \(T(n)=o(n^2)\); a frequently used stronger interpretation asks for \(T(n)=O(n^{2-\eps})\) for some \(\eps>0\). The source does not distinguish those meanings. It suggests \(O(n\log n)\) as a possible target, motivated by an algorithm for certain strongly relatively hyperbolic groups in [Fa3]. This suggested rate is not asserted as achieved, and the question does not demand proof of an optimal possible complexity. For fixed \(g,S\), constants may depend on those choices; the excerpt does not require uniformity as \(g\) varies.

The source cites Masur--Minsky, Theorem 1.3 of [MM1], for a weak-hyperbolicity property of mapping class groups. This is retained as its stated motivation, not replaced by an assertion that \(\operatorname{Mod}_g\) is a Gromov-hyperbolic group. The exact relative or action-based notion needs the cited context.

The following discussion concerns conjugacy, a different decision problem: given words representing \(x,y\), decide whether \(hxh^{-1}=y\) for some \(h\). The source calls it harder, says Hemion's original algorithm [He] appeared to give no reasonable, even exponential, time bound, proposes biautomaticity as a refinement, and says even a biautomatic structure gives only an exponential conjugacy algorithm. These are historical comments with their source scope; neither current conjugacy complexity nor biautomaticity is asserted as resolved. The mentioned following approach to conjugacy is not included in the record.
\subsection{Short English statement}
Find a subquadratic algorithm deciding whether a word in a fixed mapping class group is trivial. The source suggests \(n\log n\), while leaving the complexity model and meaning of subquadratic unspecified; its separate conjugacy discussion is preserved as context.
\subsection{Sources}
\begin{itemize}
\item Local UnsolvedMath registry, permanent ID 11000022, source code catalogue entry, Question 3.1, ``Fast word problem,'' PDF page 31. No individual proposer is supplied.
\item Farb, \emph{Problems on Mapping Class Groups and Related Topics}. Supplied URL: \url{https://www.math.uchicago.edu/~farb/papers/mcgbook.pdf}; recorded access 2026-07-29; extraction pdf-font-structure. The edited volume contains 329 explicitly typeset headings across contributed chapters, with repeated local numbering disambiguated by page.
\item Internal references [Fa3] for the strongly relatively hyperbolic comparison algorithm, Masur--Minsky [MM1], Theorem 1.3, for weak hyperbolicity, and Hemion [He] for the conjugacy algorithm. Full bibliographic entries are not supplied.
\item Supplied status partially\_solved and PARTIAL-PROGRESS. The summary reports polynomial/quadratic word algorithms and fast practical algorithms, says faster rates and curve-complex geodesic questions remain open, and frames ``fastest possible'' complexity as unresolved. The original question asks existence of a subquadratic algorithm, not absolute optimality or the curve-complex geodesic problem. These are supplied assessments without a current primary-literature audit. Background status and redistribution terms remain NEEDS\_REVIEW. The registry assigns L4 whereas background records default L3.
\end{itemize}
\end{document}
